% Section 3: Evaluation orders and the canonical linear extension.

\subsection{Plane trees, corner order, dependency order}

\begin{definition}[Plane binary tree]\label{def:plane-binary}
A \emph{plane binary tree} is a rooted tree in which every internal vertex has a left
and a right child, either of which may be a leaf; the children are ordered.  Write
$I(t)$ for the set of internal vertices and $n(t)=|I(t)|$.  A vertex is a leaf when
it has no children; every other vertex is internal and has a left and a right child
position, either of which may lead to a leaf.  Trees are drawn with the root at the top and the left
child to the left.
\end{definition}

The combinatorics of the carousel is stated for this class; the phrase \emph{plane
planted complete tree} used in the source is the same object once the root edge is
contracted, and this paper keeps the plane binary presentation throughout.

\begin{definition}[Corner order]\label{def:corner-order}
The \emph{corner order} $P(t)$ is the preorder of $I(t)$: the root, then the left
subtree, then the right subtree.  It is the left standard traversal of the source,
and it is the order in which the corners of the tree are met when the boundary is
walked from the left.
\end{definition}

\begin{definition}[Dependency order]\label{def:dependency-order}
The \emph{dependency order} on $I(t)$ is the descendant order: $u\preceq v$ when $u$
lies in the subtree of $v$.  A \emph{linear extension} of $t$ is a total order $E$ of
$I(t)$ in which every descendant precedes its ancestors; $\mathcal{L}(t)$ denotes the
set of linear extensions.
\end{definition}

A linear extension is exactly an evaluation schedule of the tree read as a program:
children must be constructed before their parent, and independent subtrees may be
interleaved arbitrarily.  This is the special case of Definition~\ref{def:schedule}
in which the program body is a tree.

\subsection{The canonical order}

\begin{definition}[Left-first postorder]\label{def:postorder}
The \emph{left-first postorder} $Q(t)$ is the order: left subtree, then right
subtree, then root.  Equivalently, traverse the tree depth-first, always visiting the
left child before the right, and list an internal vertex after both of its subtrees.
\end{definition}

\begin{definition}[Inversion count]\label{def:inversions}
For $E\in\mathcal{L}(t)$ and the corner order $P(t)$,
\[
\inv(E,P(t))=\#\{(u,v)\;:\;u\text{ precedes }v\text{ in }P(t)\ \text{and}\ v
\text{ precedes }u\text{ in }E\}.
\]
It counts the pairs on which the evaluation order departs from the corner order; it
is the Kendall tau distance between the two orders restricted to pairs that the
corner order orients.
\end{definition}

\begin{theorem}[Canonical evaluation order]\label{thm:canonical-order}
For every plane binary tree $t$,
\[
Q(t)=\Argmin_{E\in\mathcal{L}(t)}\inv(E,P(t)),
\]
and the minimizer is unique.
\end{theorem}

\begin{proof}
Induction on $n(t)$.  For $n(t)=1$ the only extension is the root and the statement
is immediate.  Let $t=\mathrm{Node}(a,b)$ with $m=n(t)$, so that
$I(t)=\{t\}\sqcup I(a)\sqcup I(b)$ and
\[
P(t)=[t]+P(a)+P(b),\qquad Q(t)=Q(a)+Q(b)+[t].
\]
Every $E\in\mathcal{L}(t)$ has the form $E=M+[t]$, where $M$ is an interleaving of
some $E_a\in\mathcal{L}(a)$ and $E_b\in\mathcal{L}(b)$: the root is last because it
depends on both subtrees, and the restrictions of $E$ to the two subtrees are linear
extensions of them.  Decompose
\[
\inv(E,P(t))=(m-1)+\inv(E_a,P(a))+\inv(E_b,P(b))+\mathrm{cross}(M),
\]
where the first term counts the root against every other vertex, and
$\mathrm{cross}(M)$ counts the inversions between a vertex of $a$ and a vertex of
$b$: pairs that $P(t)$ orders $a$-before-$b$ but $M$ orders $b$-before-$a$.

Two observations finish the proof.  First, $\mathrm{cross}(M)\ge 0$, with equality
exactly when every vertex of $a$ precedes every vertex of $b$ in $M$, that is, when
$M=E_a+E_b$ is a block.  Second, by the induction hypothesis
$\inv(E_a,P(a))$ is minimized uniquely at $Q(a)$ and $\inv(E_b,P(b))$ uniquely at
$Q(b)$.  Hence the right-hand side is minimized uniquely at
$E_a=Q(a)$, $E_b=Q(b)$, and the block interleaving, that is, at
$Q(a)+Q(b)+[t]=Q(t)$.
\end{proof}

\begin{corollary}[The canonical schedule]\label{cor:canonical-schedule}
Every plane binary tree has exactly one evaluation schedule that is closest to the
corner order in the inversion metric, and it is the left-first postorder.  In
particular the choice of schedule across incomparable branches is not free once the
criterion is fixed: it is forced.
\end{corollary}

\begin{remark}[Why the extremal criterion is the right one]\label{rem:why-criterion}
The corner order is the order in which the tree's boundary is walked; it is the only
order the planar structure supplies without a choice.  The canonical schedule is the
linear extension that departs from it as little as the dependency order allows.  The
criterion is therefore not an optimisation chosen for convenience but the closest
linearization of a planar order under a poset constraint; other criteria --- for
example preserving the left-to-right order of leaves --- are recorded as open
directions in \S\ref{sec:boundaries} because they have not been shown to select the
same extension.
\end{remark}

\subsection{Finite calibration and the limits of the statement}

\begin{calibration}[Canonical order, small trees]\label{cal:canonical-order}
For every plane binary tree with $3\le n(t)\le 6$ internal vertices, equivalently
$4\le \ell(t)\le 7$ leaves (193 trees), the left-first postorder is
the unique linear extension minimizing $\inv(\cdot,P(t))$.  The check is stored at
\path{research/iota-carousel/verify_canonical_choice.py}; the closed form of the
calibration is the statement of Theorem~\ref{thm:canonical-order}, so the calibration
is evidence for the implementation, not for the theorem.
\end{calibration}

Theorem~\ref{thm:canonical-order} is a statement about plane binary trees and their
descendant order.  It does not by itself say that an arbitrary program's dependency
order is tree-shaped, nor that a general program has a unique canonical schedule: a
program DAG may admit incomparable events that no planar structure orders.  What the
theorem does supply is the canonical choice on the planar layer.  The carousel's index
is a different object: it is the \emph{preorder} position supplied by the corner order
itself, not the canonical postorder schedule; \S\ref{sec:carousel} states the relation
between the two, and the relation is not an equality of orders.

\subsection{The planar and poset layers}

Two orders must be kept apart when this result is used.

\begin{itemize}
\item The \emph{corner order} $P(t)$ is planar: it records left and right.
\item The \emph{down-set lattice} $O(Q)$ of a dependency poset $Q$ records only
  dependency and forgets left and right; it is the lattice of Paper~I, and its
  elements are down-sets rather than orderings of vertices.
\end{itemize}

The canonical order is a statement in the planar layer.  Donaghey-type transforms act
on the planar layer as well, and the question whether they descend to $O(Q)$ has a
negative answer for the transform $\tau$ studied in the research notes
(\texttt{OQ-098}); the exact statement is not needed here and is recorded in
\S\ref{sec:boundaries}.  The distinction matters for the equivalence of
\S\ref{sec:equivalence}: a dictionary that identified the two layers would be wrong,
and the equivalence is therefore stated with the planar data retained.
